\section{Introduction and classification}
\label{sec:introduction}

The rational periodic set of a polynomial automorphism can be finite
without admitting a usable description as the coefficients vary. We
consider the complete classification problem for the family
\begin{equation}
 H_{a,b}(x,y)=(y,y^2+by+a-x),\qquad a,b\in\Z.
 \label{eq:family}
\end{equation}
The determinant of its Jacobian is \(+1\). Both the sign of the final
term and the condition that the quadratic coefficient is exactly one
are part of the problem. For a map \(H\), write
\[
 \Per(H,\Q)=\{P\in\Q^2:H^n(P)=P\text{ for some }n\ge1\}.
\]
The exact period of \(P\) is the least such positive integer \(n\).
All times in this article refer to ordinary iteration of \(H\).

A finite coordinate word \((v_0,\ldots,v_{d-1})\) denotes the orbit
whose successive points are
\((v_i,v_{i+1})\), with subscripts read modulo \(d\). Words are
identified under cyclic rotation, but not under an arbitrary permutation
of their entries. The word describes an orbit of \(H_{A,e}\) precisely
when
\begin{equation}
 v_{i-1}+v_{i+1}=v_i^2+ev_i+A
 \quad\text{for every }i.
 \label{eq:normalized_word}
\end{equation}
Repeated coordinates are allowed; the least period of the resulting
orbit must still be checked.

\begin{theorem}[Complete rational periodic set]
\label{thm:classification}
Let \(a,b\in\Z\). Write uniquely \(b=2q+e\), where \(q\in\Z\) and
\(e\in\{0,1\}\), and put
\begin{equation}
 A=a-q^2+(2-e)q.
 \label{eq:normalized_parameter}
\end{equation}
Every rational periodic point of \(H_{a,b}\) is integral. The complete
list of its periodic orbits is obtained from Table~\ref{tab:even} when
\(e=0\), or Table~\ref{tab:odd} when \(e=1\), by subtracting \(q\)
from every coordinate in the listed normalized words. At a parameter
appearing in several rows, take the union of those rows, counting
coincident orbits only once.

Every exact period belongs to \(\{1,2,3,4\}\), and
\begin{equation}
 \#\Per(H_{a,b},\Q)\le8.
 \label{eq:eight_bound}
\end{equation}
Equality holds if and only if \(e=1\) and \(A=-4\), equivalently
\(b=2q+1\) and \(a-q^2+q=-4\).
\end{theorem}

\begin{table}[htbp]
\centering
\small
\begin{tabular}{@{}L{0.22\textwidth}L{0.39\textwidth}L{0.31\textwidth}@{}}
\toprule
Parameter \(A\) & Coordinate words & Exact period and range \\
\midrule
\(1-k^2\) & \((1-k)\), \((1+k)\) & 1; \(k\ge0\); coincide at \(k=0\) \\
\(-k^2-3\) & \((-k-1,k-1)\) & 2; \(k\ge1\) \\
\(-k^2-1\) & \((-k-1,k,k)\), \((k-1,-k,-k)\) & 3; \(k\ge0\); coincide at \(k=0\) \\
\(-k^2\) & \((-k,-k,k,k)\) & 4; \(k\ge1\) \\
\bottomrule
\end{tabular}
\caption{All rational periodic orbits of \(H_{A,0}\), written as cyclic
coordinate words. The two coincidence clauses apply only to the two
words in their own row. At overlapping parameters the table means their
union; all ranges are part of the classification.}
\label{tab:even}
\end{table}

\begin{table}[htbp]
\centering
\small
\begin{tabular}{@{}L{0.25\textwidth}L{0.38\textwidth}L{0.29\textwidth}@{}}
\toprule
Parameter \(A\) & Coordinate words & Exact period and range \\
\midrule
\(-k(k+1)\) & \((-k)\), \((k+1)\) & 1; \(k\ge0\); always distinct \\
\(-k(k+1)-4\) & \((-k-2,k-1)\) & 2; \(k\ge0\) \\
\(-k(k+1)-2\) & \((-k-2,k,k)\) & 3; \(k\ge0\) \\
\(-k(k+1)-2\) & \((k-1,-k-1,-k-1)\) & 3; \(k\ge1\) \\
\(-k(k+1)-1\) & \((-k-1,-k-1,k,k)\) & 4; \(k\ge0\) \\
\bottomrule
\end{tabular}
\caption{All rational periodic orbits of \(H_{A,1}\). The second
three-cycle row starts at \(k=1\): its word at \(k=0\) would be the
fixed point \((-1,-1)\), already present in the fixed-point row at
parameter \(A=-2\). Overlapping rows again mean their union.}
\label{tab:odd}
\end{table}

The coefficient normalization, the shared symbol argument, and the
exceptional parameters are all needed for this statement. A census at
finitely many parameters would not exclude cycles at new parameters or
of larger period. Conversely, a uniform period bound would not determine
which points occur at each parameter. Here the large-parameter proof
leaves an explicitly bounded finite complement, rather than treating a
sample as evidence for the general case.

\paragraph{Relation to arithmetic finiteness and period bounds.}
The height-theoretic study of H\'enon maps is classical; foundational
work includes Silverman~\cite{silverman1994geometric}, and Ingram
develops canonical-height and specialization results
in~\cite{ingram2014canonical}. Our proof uses a direct local maximum
argument to establish integrality and does not claim a new general
finiteness theorem. For arbitrary polynomial maps of \(\Z^2\) with
integer coefficients, Pezda~\cite[Theorem 2.1]{pezda2002cycles}
determines the possible cycle lengths, with largest value 24. The task
here is different: it asks for the exact point sets in a fixed quadratic
automorphism family at every parameter.

\paragraph{Degree, coefficients, and determinant sign.}
Ingram's detailed conjecture for quadratic maps concerns
\((x,y)\mapsto(y,x+y^2+c)\) with \(c\in\Q\), whose determinant is
\(-1\); see~\cite{ingram2014canonical}. The same sign distinction is
explicit in the number-field questions of
\cite[Section 11]{dehenon2024problems}. Theorem~\ref{thm:classification}
does not resolve that conjecture, nor a uniform-boundedness conjecture
for all rational-coefficient H\'enon maps. In contrast, the constructions
of Kim, Krieger, Postolache, and Szeto~\cite{kim2025many} do have
determinant \(+1\). Their long cycles arise in growing odd degree from
rational-coefficient integer-valued polynomials, rather than from the
monic quadratic integral-coefficient family~\eqref{eq:family}.

These comparisons concern the precise objects and statements just
described, not a claim of global priority. In particular, our access to
Silverman~\cite{silverman1994geometric} was limited to its publisher
record; its arithmetic context is also discussed in the accessible
introductions of~\cite{ingram2014canonical,kim2025many}. We make no
assertion about examples or classification statements in its unread
subscription text. The proof below is self-contained.

The argument proceeds through integrality and integer translations in
Section~\ref{sec:normalization}, followed by a common six-symbol
reduction in Section~\ref{sec:six_symbols}. Section~\ref{sec:local}
classifies the resulting local equations. Section~\ref{sec:finite}
and Appendix~\ref{app:certificates} close the finite complement.
Section~\ref{sec:sharp} proves the sharp bound, and
Section~\ref{sec:returns} records the ordinary-time return law and its
scope.
